%% LyX 2.4.4 created this file.  For more info, see https://www.lyx.org/.
%% Do not edit unless you really know what you are doing.
\documentclass[english]{article}
\usepackage[T1]{fontenc}
\usepackage[latin9]{inputenc}
\usepackage{units}
\usepackage{enumitem}
\usepackage{amsmath}
\usepackage{amsthm}
\usepackage{geometry}
\geometry{verbose,tmargin=1in,bmargin=1in,lmargin=1in,rmargin=1in}

\makeatletter

%%%%%%%%%%%%%%%%%%%%%%%%%%%%%% LyX specific LaTeX commands.
\newcommand{\lyxmathsym}[1]{\ifmmode\begingroup\def\b@ld{bold}
  \text{\ifx\math@version\b@ld\bfseries\fi#1}\endgroup\else#1\fi}


%%%%%%%%%%%%%%%%%%%%%%%%%%%%%% Textclass specific LaTeX commands.
\numberwithin{equation}{section}
\numberwithin{figure}{section}
\newlength{\lyxlabelwidth}      % auxiliary length 

%%%%%%%%%%%%%%%%%%%%%%%%%%%%%% User specified LaTeX commands.
\usepackage{fancyhdr}
\usepackage{lastpage}
\pagestyle{fancy}
\usepackage{enumitem}
\usepackage{ifthen}
\everymath=\expandafter{\the\everymath\displaystyle}
%\DeclareMathSizes{10}{10}{10}{10}
\usepackage{multicol}

\usepackage{tikz}
\usepackage{tikz-3dplot}
\usetikzlibrary{calc,arrows}

\usetikzlibrary{patterns}
\usetikzlibrary{plotmarks}
\usepackage{pgfplots}
\pgfplotsset{compat=newest}
\usepgfplotslibrary{polar}

\usepackage{calc}
\usepackage[nomessages]{fp}

\usepackage{datetime}
% Added by lyx2lyx
\setlength{\parskip}{\smallskipamount}
\setlength{\parindent}{0pt}

\makeatother

\usepackage{babel}
\begin{document}
\renewcommand{\author}{Dr.\,James Ashe}

\newcommand{\classNum}{331}

\newcommand{\classSec}{A}

\newcommand{\examType}{Exam 3}

\renewcommand{\date}{\formatdate{4}{11}{2013}}

%%First page's header%%%%%%%%%%%%%%%%%%%%%%
\fancypagestyle{firststyle} %%header settings for first pagr
{
	\fancyhead[L]{MTH \classNum{}(\classSec{}) \author{}\\
	\textbf{\examType{}} ~\date{}}
	\fancyhead[C]{\textbf{Name:}}
	%\fancyhead[R]{\textbf{Name:\hspace{5cm}}}
	\setlength{\headsep}{14pt}
}
%%%%%%%%%%%%%%%%%%%%%%%%%%%%%%%%%%%%%%%%%%

%%Next page's header%%%%%%%%%%%%%%%%%%%%%%%
\setlist{nolistsep}
\lhead{MTH \classNum{}(\classSec{})} 
\chead{\textbf{Name:}} 
\cfoot{\thepage{}~of~\pageref{LastPage}} 
\setlength{\headsep}{5pt} 
%%%%%%%%%%%%%%%%%%%%%%%%%%%%%%%%%%%%%%%%%%%

%%Various settings; see the preamble for other settings%%%%%
%\setenumerate[0]{leftmargin=12pt,itemindent=0pt,labelwidth=0pt}
\setlist{topsep=.5pt}
\renewcommand{\labelenumi}{\textbf{\arabic{enumi}.}}
\renewcommand{\labelenumii}{\textbf{\alph{enumii}.}}
\thispagestyle{firststyle}
%%%%%%%%%%%%%%%%%%%%%%%%%%%%%%%%%%%%%%%%%%

\newcommand{\xmax}{0}
\newcommand{\xmin}{-\xmax}
\newcommand{\xstep}{0}
\newcommand{\ymax}{0}
\newcommand{\ymin}{-\ymax}
\newcommand{\ystep}{0} 
\newcommand{\dspace}{\vspace{1cm}}
\newcommand{\xa}{0}
\newcommand{\xb}{0}
\newcommand{\ya}{1}
\newcommand{\yb}{1}

\newcommand{\xspace}{\vspace{1cm}}

\textbf{Unless otherwise indicated, each question is of equal value
{[}4 pts{]}. Show all of your work.}

\emph{Use $\mathbf{u}=\left\langle 2,4,1\right\rangle $ and $\mathbf{v}=\left\langle 1,-3,2\right\rangle $
to perform the indicated operations in problems \ref{enu:vect-start}--\ref{enu:vect-end}.}
\begin{enumerate}[resume]
\item \label{enu:vect-start}$\mathbf{u}-2\mathbf{v}$\vfill
\item $\mathbf{u\cdot}\mathbf{v}$\vfill
\item \label{enu:vect-end}$\mathbf{u\times v}$\vfill\vspace{1cm}
\end{enumerate}
\emph{Use the graph to draw the following vectors in problems \ref{enu:graph-start}--\ref{enu:graph-end}
(provide labels). Make certain to indicate direction.}
\begin{enumerate}[resume]
\item \label{enu:graph-start}\begin{multicols}{2}

\begin{enumerate}[resume]
\item \textbf{$\nicefrac{1}{2}\mathbf{u}$}
\item $\mathbf{u}+\mathbf{v}$ 
\item $\mathbf{u}-\mathbf{v}$ (you do not have to draw it as a position
vector)
\end{enumerate}
\vfill\columnbreak

\tdplotsetmaincoords{0}{0}% %%specifes orientation %%\tdplotsetrotatedcoords{0}{20}{0}% %%specifes angle of rotation
\begin{tikzpicture} 		
	[tdplot_main_coords,
		scale=.5,
		grid/.style={very thin,gray}, 
		axis/.style={->,black,thick},
		vector/.style={very thick,red,-stealth},
		vector guide/.style={dashed,red,thick},
		angle/.style={black,thick}]
	
	%draw a grid in the x-y plane 	
	%\foreach \x in {-1,0,...,1} 		
	%\foreach \y in {-1,0,...,1} 		
	{ 			
		%\draw[grid] (\x,-0.25) -- (\x,1.25); 			
		%\draw[grid] (-0.25,\y) -- (1.25,\y); 		
	} 		
	
	
	
	%draw the axes 	
	\draw[axis] (-1,0) -- (5,0) node[anchor=west]{$x$}; 	
	\draw[axis] (0,-1) -- (0,5) node[anchor=west]{$y$}; 	

	%standard tikz coordinate definition using x, y, z coords 	
	\coordinate (O) at (0,0);

	%draw vector
	\draw[->, very thick,red,-stealth] (0,0) -- (4,4) node[anchor=south west]{\textbf{u}} node[anchor=north east,pos=.6]{$$};
	\draw[->,very thick,blue,-stealth] (0,0) -- (3,1) node[anchor=south west]{\textbf{v}} node[anchor=north east,pos=.6]{$$};
	%\draw[very thick,blue,stealth-] (1.75,.5) -- (.4,1.5) node[anchor=south west,pos=.65]{$A-B$} node[anchor=north east,pos=.6]{$c$};
	%\draw[black] (0,0) -- (0,0) node[anchor=south west]{$\theta$};
	%draw an arc illustrating the angle defining the orientation 
	%\tdplotsetthetaplanecoords{45}	
	%\tdplotdrawarc[tdplot_rotated_coords,angle]{(O)}{.35}{54.7}{90}{anchor=west}{$\phi$} %{center}{radius}{start}{end angle}
\end{tikzpicture}
	\end{multicols}
\item \label{enu:graph-end}\begin{multicols}{2}

\begin{enumerate}[resume]
\item $\mbox{proj}_{\mathbf{v}}\mathbf{u}$
\item $\mbox{proj}_{\mathbf{u}}\mathbf{v}$\vfill
\end{enumerate}
\columnbreak

\tdplotsetmaincoords{0}{0}% %%specifes orientation %%\tdplotsetrotatedcoords{0}{20}{0}% %%specifes angle of rotation
\begin{tikzpicture} 		
	[tdplot_main_coords,
		scale=.5,
		grid/.style={very thin,gray}, 
		axis/.style={->,black,thick},
		vector/.style={very thick,red,-stealth},
		vector guide/.style={dashed,red,thick},
		angle/.style={black,thick}]
	
	%draw a grid in the x-y plane 	
	%\foreach \x in {-1,0,...,1} 		
	%\foreach \y in {-1,0,...,1} 		
	{ 			
		%\draw[grid] (\x,-0.25) -- (\x,1.25); 			
		%\draw[grid] (-0.25,\y) -- (1.25,\y); 		
	} 		
	
	
	
	%draw the axes 	
	\draw[axis] (-1,0) -- (5,0) node[anchor=west]{$x$}; 	
	\draw[axis] (0,-1) -- (0,5) node[anchor=west]{$y$}; 	

	%standard tikz coordinate definition using x, y, z coords 	
	\coordinate (O) at (0,0);

	%draw vector
	\draw[->, very thick,red,-stealth] (0,0) -- (3,4) node[anchor=south west]{\textbf{u}} node[anchor=north east,pos=.6]{$$};
	\draw[->,very thick,blue,-stealth] (0,0) -- (1.5,0) node[anchor=south west]{\textbf{v}} node[anchor=north east,pos=.6]{$$};
	%\draw[very thick,blue,stealth-] (1.75,.5) -- (.4,1.5) node[anchor=south west,pos=.65]{$A-B$} node[anchor=north east,pos=.6]{$c$};
	%\draw[black] (0,0) -- (0,0) node[anchor=south west]{$\theta$};
	%draw an arc illustrating the angle defining the orientation 
	%\tdplotsetthetaplanecoords{45}	
	%\tdplotdrawarc[tdplot_rotated_coords,angle]{(O)}{.35}{54.7}{90}{anchor=west}{$\phi$} %{center}{radius}{start}{end angle}
\end{tikzpicture}
	\end{multicols}
\end{enumerate}
\emph{Use the graph to perform the indicated operations (provide labels). }
\begin{enumerate}[resume]
\item \begin{multicols}{2}

\begin{enumerate}[resume]
\item Use the right-hand rule to label the \emph{x }and \emph{y }axes. 
\item Plot the point $\left(4,2,3\right)$.
\item Sketch the plane $\left(x,y,3\right)$.
\end{enumerate}
\columnbreak\tdplotsetmaincoords{65}{105}% %%specifes orientation
%%\tdplotsetrotatedcoords{0}{20}{0}% %%specifes angle of rotation
\begin{tikzpicture} 		
	[tdplot_main_coords,
		scale=2,
		cube/.style={purple}, 
		grid/.style={very thin,gray}, 
		axis/.style={->,black,thick},
		vector/.style={very thick,red,-stealth},
		vector guide/.style={dashed,red,thick},
		angle/.style={black,thick}]
	
	%draw a grid in the x-y plane 	
	\foreach \x in {-0.25,0,...,1.25} 		
	\foreach \y in {-0.25,0,...,1.25} 		
	{ 			
		\draw[grid] (\x,-0.25) -- (\x,1.25); 			
		\draw[grid] (-0.25,\y) -- (1.25,\y); 		
	} 		
		
	%draw the axes 	
	\draw[axis] (0,0,0) -- (1.5,0,0) node[anchor=west]{$$}; 	
	\draw[axis] (0,0,0) -- (0,1.5,0) node[anchor=west]{$$}; 	
	\draw[axis] (0,0,0) -- (0,0,1.5) node[anchor=west]{$z$};
	
	
	%standard tikz coordinate definition using x, y, z coords 	
	\coordinate (O) at (0,0,0);
	%tikz-3dplot coordinate definition using r, theta, phi coords 	
	\tdplotsetcoord{P}{.8}{55}{60}

	%draw vector
	%\draw[vector] (O) -- (P);
	%\draw[vector guide] (O) -- (Pxy); %\draw[vector guide] (Pxy) -- (P);
	%\draw[->, very thick,red,-stealth] (0,0,0) -- (1,1,1);
	%\draw[->,very thick,blue,-stealth] (0,0,0) -- (1,1,0);
	%draw an arc illustrating the angle defining the orientation 
	%\tdplotsetthetaplanecoords{45}	
	%\tdplotdrawarc[tdplot_rotated_coords,angle]{(O)}{.35}{54.7}{90}{anchor=west}{$\phi$} %{center}{radius}{start}{end angle}
\end{tikzpicture}
	\end{multicols}\newpage
\end{enumerate}
\emph{Use the vectors $\mathbf{u}=\left\langle 3,6\right\rangle $,
$\mathbf{v}=\left\langle 2,3\right\rangle $ and the point $P(1,-4)$
to complete problems \ref{enu:pts_and_vects-start}--\ref{enu:pts_and_vects-end}.}
\begin{enumerate}[resume]
\item \label{enu:pts_and_vects-start}Find the line through $P$ and parallel
to $\mathbf{u}$.\xspace\xspace\vfill
\item Find the angle between $\mathbf{u}$ and $\mathbf{v}$.\xspace\xspace\vfill
\item \label{enu:Find-a-vector}Find a vector orthogonal to $\mathbf{u}$
and $\mathbf{v}$.\xspace\xspace\vfill
\item \label{enu:pts_and_vects-end}Is the vector found in problem \ref{enu:Find-a-vector}
unique? If not find another vector orthogonal to $\mathbf{u}$ and
$\mathbf{v}$.\xspace\vfill
\end{enumerate}
\emph{Let $\mathbf{r}(t)=\left\langle 1,2t^{-2},e^{-t}\right\rangle $.
Evaluate problems \ref{enu:calc}--\ref{clac-end}.}
\begin{enumerate}[resume]
\item \label{enu:calc}$\mathbf{r}'(t)$\xspace\vfill
\item ${\displaystyle \int}_{0}^{1}\mathbf{r}(t)\,dt$\xspace\xspace\vfill
\item \label{clac-end}${\displaystyle \lim_{t\rightarrow\infty}\mathbf{r}(t)}$\xspace\vfill
\end{enumerate}
\newpage

\emph{Compute the slope of the tangent line of the vector valued function
at $t=1$.}
\begin{enumerate}[resume]
\item $\mathbf{r}(t)=\left\langle 1+t^{3},2\sin3t,e^{t}\right\rangle $\vfill
\end{enumerate}
\emph{Complete problems \ref{first word problem}--\ref{enu:last word problem}.}
\begin{enumerate}[resume]
\item \label{first word problem}A sailboat floats in a current that flows
due west at 3 m/s. The wind is blowing $30\textdegree$ north of east
at a speed of 5 m/s. Find the speed of the boat relative to the shore.\xspace\vfill
\item A constant force vector $\mathbf{F}=\left\langle 50,120\right\rangle $
(pounds) is used to move a heavy object 50 feet along the \emph{x-}axis.
Calculate the work done.\xspace\vfill
\item \label{enu:last word problem}A 40-pound weight hangs from the end
of a 2-foot bar that is attached to a wall at an angle of $45^{\circ}$
below the horizontal. 

\begin{enumerate}[resume]
\item Find the magnitude of the torque about the point on the wall where
the bar is attached.\xspace\vfill
\item Draw a picture and indicate the direction of the torque.\xspace\xspace
\end{enumerate}
\end{enumerate}

\end{document}
